% =============================================================================
%  Part I: Spherical Scatterer  (Sections 1--10)
% =============================================================================

\section{Introduction}

The scattering of elastic waves by localised heterogeneities is a fundamental
problem in seismology, ultrasonics, and geophysical exploration.  When the
scatterer is a sphere of radius~$a$ embedded in a homogeneous isotropic
background, the problem admits an exact solution in terms of partial-wave (Mie)
series \citep{Korneev1993,Gritto1995}.  In the long-wavelength (Rayleigh) regime
$ka\ll 1$, only the lowest-order multipoles contribute and the scattered field
can be expressed compactly through an effective T-matrix that maps the incident
wavefield to the scattered wavefield.

The T-matrix concept originates in quantum scattering theory and was adapted to
elastic-wave problems by \citet{Waterman1976}; \citet{Bostock1991}
and \citet{BostockKennett1992} applied it to surface-wave
scattering.  \citet{Pollitz1998} derived the scattered field for a small
spherical inclusion without assuming small perturbations, arriving at
self-consistent amplification factors that are functionally equivalent to
renormalised material contrasts.  \citet{Pellegrini2011}
formalised this as a mean-field T-matrix approach: by approximating the interior
stress and momentum fields by their volume averages, the full
multiple-scattering series inside the inclusion can be summed in closed form,
yielding a T-matrix that satisfies the optical theorem exactly.

On the computational side, Jakobsen and co-workers developed integral-equation
methods for seismic modelling and inversion based on the Lippmann--Schwinger
equation.  \citet{Jakobsen2012} introduced the T-matrix formalism for
acoustic forward modelling; \citet{JakobsenUrsin2015} extended
this to direct iterative T-matrix inversion; and \citet{JakobsenWu2016}
proposed a renormalised scattering series for frequency-domain waveform modelling
in the presence of strong velocity contrasts, overcoming the convergence
limitations of the Born series.  Most recently, \citet{Shekhar2024} formulated a numerical integral-equation method for
microseismic wavefield modelling in arbitrarily anisotropic elastic media,
providing the Green's tensor expressions and the Lippmann--Schwinger framework
used as the starting point of this work.

In this note, we apply symbolic computer algebra (SymPy) to evaluate the
Lippmann--Schwinger integral equation analytically for two scatterer geometries:
a sphere and a cube, both with constant isotropic contrasts $\Delta\lambda$,
$\Delta\mu$, $\Delta\rho$.  The sphere yields closed-form results with
isotropic effective contrasts; the cube introduces cubic anisotropy in the
effective stiffness, providing the natural building block for effective medium
theories of heterogeneous solids.


% =====================================================================
\part*{Part I: Spherical Scatterer}
% =====================================================================

\section{Problem Statement}

We consider the scattering of elastic waves by a small spherical inclusion of
radius~$a$ embedded in an infinite homogeneous isotropic background medium with
P-wave velocity~$\alpha$, S-wave velocity~$\beta$, and density~$\rho$.  The
inclusion has constant isotropic contrasts $\Delta\lambda$, $\Delta\mu$,
$\Delta\rho$ relative to the background.

The wavefield satisfies the Lippmann--Schwinger integral equation (Eq.~8
of \citealt{Shekhar2024}):
\begin{equation}
  u_i(\bx) = u_i^0(\bx) + \int_D \Bigl[
    G_{ij}(\bx,\bx')\,\omega^2\,\Delta\rho\;u_j(\bx')
    + H_{ijk}(\bx,\bx')\,\Delta c_{jklm}\,\varepsilon_{lm}(\bx')
  \Bigr] d^3x',
\label{eq:LS}
\end{equation}
where the strain--Green's tensor is
\begin{equation}
  H_{ijk}^{(0)}(\bx,\bx')
  = \tfrac{1}{2}\bigl(G_{ij,k} + G_{ik,j}\bigr),
\label{eq:Hdef}
\end{equation}
and the isotropic stiffness contrast is
\begin{equation}
  \Delta c_{jklm}
  = \Delta\lambda\,\delta_{jk}\,\delta_{lm}
  + \Delta\mu\,(\delta_{jl}\,\delta_{km} + \delta_{jm}\,\delta_{kl}).
\label{eq:Dc}
\end{equation}


\section{Background Green's Tensor}

The elastodynamic Green's tensor for a homogeneous isotropic medium
(Appendix~A of \citealt{Shekhar2024}) admits the radial decomposition:
\begin{equation}
  G_{ij}(\bx) = f(r)\,\delta_{ij} + g(r)\,\hat{n}_i\,\hat{n}_j,
\label{eq:Grad}
\end{equation}
with $\hat{n}_i = x_i/r$ and radial functions
\begin{equation}
  X(r) = \frac{e^{i\omega r/\alpha}}{4\pi\rho\alpha^2}, \quad
  V(r) = \frac{e^{i\omega r/\beta}}{4\pi\rho\beta^2}, \quad
  C = \frac{1-\beta^2/\alpha^2}{8\pi\rho\beta^2},
  \qquad
  f(r) = \frac{V(r)}{r} - \frac{C}{r},
\label{eq:fdef}
\end{equation}
\begin{equation}
  g(r) = \frac{3C}{r} - \frac{X(r)}{r} + \frac{V(r)}{r} .
\label{eq:gdef}
\end{equation}


\section{Calculus: Spherical Coordinate Reduction}

The central computational challenge is to evaluate volume integrals of the
form $\int_{\text{sphere}} G_{ij}\,d^3x$,
$\int_{\text{sphere}} G_{ij,k}\,d^3x$ and
$\int_{\text{sphere}} G_{ij,kl}\,d^3x$ over a sphere of radius~$a$ centred
at the scatterer.  The key insight is that although these integrals live in
Cartesian tensor space, the radial decomposition~\eqref{eq:Grad} allows them
to be factored into angular integrals (which evaluate to known isotropic
tensor identities) multiplied by 1D radial integrals (which SymPy evaluates
in closed form).

\subsection{Conversion to Spherical Coordinates}

We write the volume element as $d^3x = r^2\,dr\,d\Omega$, where
$d\Omega = \sin\theta\,d\theta\,d\varphi$ is the solid-angle element.  Every
Cartesian integrand is expressed in terms of the radial coordinate~$r$ and
the unit direction vector $\hat{n}_i = x_i/r$.  Because $f(r)$ and $g(r)$
depend only on~$r$, the angular and radial parts separate cleanly:
\begin{equation}
  \int_{\text{sphere}}
    \bigl[f(r)\,\delta_{ij} + g(r)\,\hat{n}_i\,\hat{n}_j\bigr]
    r^2\,dr\,d\Omega
  = \delta_{ij}\int_0^a r^2 f(r)\,dr \int d\Omega
  + \int_0^a r^2 g(r)\,dr \int d\Omega\;\hat{n}_i\hat{n}_j.
\label{eq:sphfact}
\end{equation}

\subsection{Angular Integral Identities}

The angular integrals of products of direction cosines over the unit sphere
are determined entirely by symmetry.  The only non-vanishing cases with an
even number of $\hat{n}$ factors are:
\begin{equation}
  \int d\Omega = 4\pi,
\label{eq:ang0}
\end{equation}
\begin{equation}
  \int d\Omega\;\hat{n}_i\,\hat{n}_j = \frac{4\pi}{3}\,\delta_{ij},
\label{eq:ang2}
\end{equation}
\begin{equation}
  \int d\Omega\;\hat{n}_i\,\hat{n}_j\,\hat{n}_k\,\hat{n}_l
  = \frac{4\pi}{15}\,(\delta_{ij}\delta_{kl}
    + \delta_{ik}\delta_{jl} + \delta_{il}\delta_{jk}).
\label{eq:ang4}
\end{equation}
Integrals with an odd number of $\hat{n}$ factors vanish by parity.  This is
the reason that $\int H_{ijk}\,d^3x = 0$: each term in $G_{ij,k}$ contains
an odd number of direction cosines after the chain rule is applied, so the
angular integral is zero.

\subsection{Chain-Rule Decomposition of Derivatives}

The Cartesian derivatives of $G_{ij}$ must be expressed in terms of
$\hat{n}_i$ and radial derivatives of $f$ and~$g$ before integration.  We
use the fundamental identities:
\begin{equation}
  \frac{\partial r}{\partial x_k} = \hat{n}_k, \qquad
  \frac{\partial\hat{n}_i}{\partial x_k}
  = \frac{1}{r}\bigl(\delta_{ik} - \hat{n}_i\,\hat{n}_k\bigr).
\label{eq:chain}
\end{equation}

\textbf{First derivative.}  Applying the chain rule to
$G_{ij} = f\,\delta_{ij} + g\,\hat{n}_i\hat{n}_j$:
\begin{equation}
  G_{ij,k} = f'\,\delta_{ij}\,\hat{n}_k
  + g'\,\hat{n}_i\,\hat{n}_j\,\hat{n}_k
  + \frac{g}{r}\bigl(\delta_{ik}\,\hat{n}_j
    + \delta_{jk}\,\hat{n}_i - 2\hat{n}_i\,\hat{n}_j\,\hat{n}_k\bigr),
\label{eq:Gijk}
\end{equation}
where primes denote $d/dr$.  Every term contains an odd number of
$\hat{n}$ factors (one or three), confirming that
$\int G_{ij,k}\,d^3x = 0$.

\textbf{Second derivative.}  Differentiating~\eqref{eq:Gijk} with respect
to~$x_l$ and again applying~\eqref{eq:chain} produces terms with angular
structures involving 0, 2, or 4 direction cosines:
\begin{equation}
  G_{ij,kl} = \text{terms proportional to}
  \;\delta\delta,\;\delta\hat{n}\hat{n},\;\hat{n}\hat{n}\hat{n}\hat{n},
\label{eq:Gijkl}
\end{equation}
with coefficients built from $f$, $g$ and their first and second radial
derivatives.  Rather than writing out all terms explicitly, we use a
\emph{contraction strategy}: the full tensor
$\int G_{ij,kl}\,d^3x$ must have the isotropic form
$A\,\delta_{ij}\delta_{kl} + B\,(\delta_{ik}\delta_{jl}+\delta_{il}\delta_{jk})$,
and we can determine the two unknowns $A$ and~$B$ from any two independent
contractions.

\subsection{Contraction Strategy for $A$ and $B$}

\textbf{Contraction~1:}  Trace over $(i,j)$.  Setting $i=j$ and summing
gives $G_{ii} = 3f + g \equiv h(r)$, which is a scalar function of~$r$.
Its second Cartesian derivative is simply:
\begin{equation}
  G_{ii,kl} = h''(r)\,\hat{n}_k\,\hat{n}_l
  + \frac{h'(r)}{r}\bigl(\delta_{kl} - \hat{n}_k\,\hat{n}_l\bigr).
\label{eq:trace1}
\end{equation}
Integrating over the sphere using \eqref{eq:ang0}--\eqref{eq:ang2}:
\begin{equation}
  \int G_{ii,kl}\,d^3x
  = \delta_{kl} \int_0^a r^2
    \left[\frac{4\pi}{3}\,h'' + \frac{8\pi}{3}\,\frac{h'}{r}\right] dr
  = (3A + 2B)\,\delta_{kl}.
\label{eq:contr1}
\end{equation}

\textbf{Contraction~2:}  Divergence structure.  Summing over $j=k$ gives
the divergence $G_{ij,j} = q(r)\,\hat{n}_i$, where
$q(r) = f' + g' + 2g/r$.  Its further derivative with respect to~$x_l$ has
the same scalar-times-vector structure:
\begin{equation}
  G_{ij,jl} = q'(r)\,\hat{n}_i\,\hat{n}_l
  + \frac{q(r)}{r}\bigl(\delta_{il} - \hat{n}_i\,\hat{n}_l\bigr).
\label{eq:divG}
\end{equation}
Integrating:
\begin{equation}
  \int G_{ij,jl}\,d^3x
  = \delta_{il} \int_0^a r^2
    \left[\frac{4\pi}{3}\,q' + \frac{8\pi}{3}\,\frac{q}{r}\right] dr
  = (A + 4B)\,\delta_{il}.
\label{eq:contr2}
\end{equation}
Solving the $2\times 2$ system \eqref{eq:contr1}--\eqref{eq:contr2} yields
$A$ and~$B$, with only 1D radial integrals for SymPy to evaluate.


\section{Analytical Volume Integrals}

\subsection{Green's Tensor Integral}

Using the decomposition of Section~4, the volume integral of the Green's
tensor reduces to:
\begin{equation}
  \int_{\text{sphere}} G_{ij}\,d^3x
  = \delta_{ij} \int_0^a r^2
    \left[4\pi\,f(r) + \frac{4\pi}{3}\,g(r)\right] dr
  \equiv \Gamma_0\,\delta_{ij}.
\label{eq:Gamma0}
\end{equation}
After exact radial evaluation by SymPy:
\begin{equation}
  \Gamma_0 = \frac{%
    -3\alpha\beta
    + 2\alpha\bigl(\beta - i\omega a\bigr)\,e^{i\omega a/\beta}
    + \beta\bigl(\alpha - i\omega a\bigr)\,e^{i\omega a/\alpha}
  }{3\alpha\beta\,\omega^2\rho}.
\label{eq:Gamma0result}
\end{equation}

\subsection{Second-Derivative Integral}

As derived in Section~4.4, the volume integral of $G_{ij,kl}$ has the
isotropic structure
\begin{equation}
  \int_{\text{sphere}} G_{ij,kl}(\bx)\,d^3x
  = A\,\delta_{ij}\delta_{kl}
  + B\,(\delta_{ik}\delta_{jl} + \delta_{il}\delta_{jk}),
\label{eq:ABiso}
\end{equation}
with $A$ and $B$ determined by the two contractions
\eqref{eq:contr1}--\eqref{eq:contr2}.  Defining
$\xi_P = \omega a/\alpha$ and $\xi_S = \omega a/\beta$, the exact
closed-form results are:
\begin{equation}
  A = \frac{%
    2\alpha^3\beta\bigl(1-e^{i\xi_S}\bigr)
    + 4i\alpha^3\omega a\,e^{i\xi_S}
    + 3\alpha\beta^3\bigl(1-e^{i\xi_P}\bigr)
    + i\beta^3\omega a\,e^{i\xi_P}
  }{15\alpha^3\beta^3\rho},
\label{eq:Aresult}
\end{equation}
\begin{equation}
  B = \frac{%
    2\alpha^3\beta\bigl(1-e^{i\xi_S}\bigr)
    - i\alpha^3\omega a\,e^{i\xi_S}
    - 2\alpha\beta^3\bigl(1-e^{i\xi_P}\bigr)
    + i\beta^3\omega a\,e^{i\xi_P}
  }{15\alpha^3\beta^3\rho}.
\label{eq:Bresult}
\end{equation}

\smallskip
\noindent
\textit{Note:}  The parity argument shows $\int H_{ijk}\,d^3x = 0$, so the
stiffness contrast enters only through the strain equation.


\section{T-Matrix Coupling Coefficients}

The T-matrix tensor, defined by contracting the symmetrised
second-derivative integral with the stiffness contrast, is isotropic:
\begin{equation}
  T_{mnlp} = T_1\,\delta_{mn}\delta_{lp}
  + T_2\,(\delta_{ml}\delta_{np} + \delta_{mp}\delta_{nl}),
\label{eq:Tiso}
\end{equation}
where:
\begin{multline}
  T_1 = \frac{1}{15\alpha^3\beta^3\rho}\Bigl[
    \Delta\lambda\Bigl(
      10\alpha^3\beta\bigl(1-e^{i\xi_S}\bigr)
      + 5\alpha\beta^3\bigl(e^{i\xi_P}-1\bigr)
      + 5i\beta^3\omega a\,e^{i\xi_S}
    \Bigr) \\
    + \Delta\mu\Bigl(
      4\alpha^3\beta\bigl(1-e^{i\xi_S}\bigr)
      - 2i\alpha^3\omega a\,e^{i\xi_S}
      + 4\alpha\beta^3\bigl(e^{i\xi_P}-1\bigr)
      + 2i\beta^3\omega a\,e^{i\xi_P}
    \Bigr)
  \Bigr],
\label{eq:T1}
\end{multline}
with $\xi_P = \omega a/\alpha$ and $\xi_S = \omega a/\beta$.
\begin{equation}
  T_2 = \frac{\Delta\mu}{15\alpha^3\beta^3\rho}\Bigl[
    4\alpha^3\beta\bigl(1-e^{i\xi_S}\bigr)
    + 3i\alpha^3\omega a\,e^{i\xi_S}
    + \alpha\beta^3\bigl(1-e^{i\xi_P}\bigr)
    + 2i\beta^3\omega a\,e^{i\xi_P}
  \Bigr].
\label{eq:T2}
\end{equation}


\section{Self-Consistent Solution}

Under the assumption that the wavefield is approximately constant inside
the small sphere, the Lippmann--Schwinger equation reduces to a
self-consistent algebraic system that decouples into three independent
scalar equations for displacement, dilatation, and deviatoric strain, with
amplification factors:
\begin{equation}
  A_u = \frac{1}{1-\omega^2\Delta\rho\,\Gamma_0}, \qquad
  A_\theta = \frac{1}{1-3T_1-2T_2}, \qquad
  A_e = \frac{1}{1-2T_2}.
\label{eq:amplsph}
\end{equation}


\section{Effective (Renormalised) Contrasts --- The T-Matrix}

The scattered field takes the Born form with renormalised contrasts:
\begin{equation}
  u_i^{\text{scat}}(\bx)
  = V\Bigl[
    H_{ijk}(\bx,\bx_s)\,\bigl(
      \Delta\lambda^*\,\theta^0\,\delta_{jk}
      + 2\Delta\mu^*\,\varepsilon_{jk}^0\bigr)
    + \omega^2\,G_{ij}(\bx,\bx_s)\,\Delta\rho^*\,u_j^0
  \Bigr],
\label{eq:scatBorn}
\end{equation}
where $V = \tfrac{4}{3}\pi a^3$ is the sphere volume.  The effective
contrasts are:

\medskip
\noindent\textbf{Density:}
\begin{equation}
  \Delta\rho^* = \frac{\Delta\rho}{1-\omega^2\Delta\rho\,\Gamma_0}.
\label{eq:Drhostar}
\end{equation}

\noindent\textbf{Shear modulus:}
\begin{equation}
  \Delta\mu^* = \frac{\Delta\mu}{1-2T_2}.
\label{eq:Dmustar}
\end{equation}

\noindent\textbf{Lam\'e parameter:}
\begin{equation}
  \Delta\lambda^* = \frac{\Delta\lambda + \tfrac{2}{3}\Delta\mu}
  {1 - 3T_1 - 2T_2}
  - \frac{\tfrac{2}{3}\Delta\mu}{1 - 2T_2}.
\label{eq:Dlamstar}
\end{equation}


\section{Born Approximation Verification}

In the Born (point-scatterer) limit $a\to 0$, all amplification factors
tend to unity:
\begin{equation}
  \lim_{a\to 0}\Delta\rho^* = \Delta\rho, \qquad
  \lim_{a\to 0}\Delta\lambda^* = \Delta\lambda, \qquad
  \lim_{a\to 0}\Delta\mu^* = \Delta\mu.
\label{eq:Bornlim}
\end{equation}
All three limits have been verified analytically with SymPy.


\section{Numerical Example}

Background: $V_P = 5000$~m/s, $V_S = 3000$~m/s,
$\rho = 2500$~kg/m$^3$.  Contrasts: $\Delta\lambda = 2$~GPa,
$\Delta\mu = 1$~GPa, $\Delta\rho = 100$~kg/m$^3$.
Frequency $f = 10$~Hz, sphere radius $a = 10$~m.

\medskip
\begin{center}
\begin{tabular}{ll}
\toprule
\textbf{Quantity} & \textbf{Value} \\
\midrule
$k_P a$ & $0.1257$ \\
$k_S a$ & $0.2094$ \\
\midrule
$\Gamma_0$ & $1.730890\mathrm{e}{-09} + 2.282529\mathrm{e}{-10}\,i$ \\
$A$        & $-3.781591\mathrm{e}{-13} + 9.278517\mathrm{e}{-13}\,i$ \\
$B$        & $2.248959\mathrm{e}{-13} - 1.438708\mathrm{e}{-12}\,i$ \\
$T_1$      & $1.492641\mathrm{e}{-03} - 1.253138\mathrm{e}{-02}\,i$ \\
$T_2$      & $-1.532633\mathrm{e}{-04} - 5.108564\mathrm{e}{-04}\,i$ \\
$A_u$      & $1.000684\mathrm{e}{+00} + 9.023389\mathrm{e}{-05}\,i$ \\
$A_\theta$ & $1.002681\mathrm{e}{+00} - 3.888157\mathrm{e}{-02}\,i$ \\
$A_e$      & $9.996925\mathrm{e}{-01} - 1.021086\mathrm{e}{-03}\,i$ \\
$\Delta\rho^*$     & $1.000684\mathrm{e}{+02} + 9.023389\mathrm{e}{-03}\,i$ \\
$\Delta\lambda^*$  & $2.007355\mathrm{e}{+09} - 1.030035\mathrm{e}{+08}\,i$ \\
$\Delta\mu^*$      & $9.996925\mathrm{e}{+08} - 1.021086\mathrm{e}{+06}\,i$ \\
\bottomrule
\end{tabular}
\end{center}
